\section{Conclusion}
\label{sec:conclusion}

We presented a knowledge-based belief-space planning framework that integrates proactive human querying into robot decision making. The robot maintains uncertainty over symbolic possible worlds, uses POMCP for online planning, and requests human feedback only when the observation-updated belief remains insufficiently confident. By selecting queries that are expected to reduce entropy over task-relevant frontier states, the system decides when and what to ask after ambiguous observations. This formulation uses human feedback as a belief-refinement mechanism rather than as action-level control, enabling robots to reduce unnecessary supervision while resolving ambiguity before it causes task failure. The user study further evaluates whether this selective state-level interaction reduces subjective workload compared with continuous supervision or action-level assistance.
